\appendix

\section{Cohort Construction: Full Derivation} \label{app:cohort}

KFRE's inputs are drawn from separate, asynchronously-ordered labs (serum creatinine, a urine albumin/creatinine test, and, for the 8-variable form, a chemistry panel) — unlike a design that avoids multi-lab merging altogether by encoding one row per single-lab event with missingness flags (the twenty\_features scenario's approach, Section~\ref{sec:cohort}). Producing complete covariate rows from asynchronously measured inputs for four\_features/eight\_features required a merge design with two design choices: (1) whether the row format should still be time-varying given that KFRE itself is a point-in-time equation, and (2) how to handle a patient having multiple draws of the same lab within one admission.

\textbf{Row format.} We keep the time-varying (\texttt{start}, \texttt{stop}) format used throughout. KFRE's formula is agnostic to row packaging — it can be applied independently to each row, treating every row as its own landmark snapshot in the sense of landmarking methodology for time-varying covariates in survival analysis \citep{vanhouwelingen2007landmarking,putter2022landmarking2}.

\textbf{Merge rule: creatinine-anchored nearest-value matching.} Each row is anchored on one creatinine measurement (eGFR is computed from it, via the CKD-EPI 2021 equation, at that instant). Other required labs are attached to the anchor by nearest-value matching using the matching rules below rather than by unbounded last-observation-carried-forward, which is known to bias sparse EHR-derived time-varying covariates \citep{hu2022locf_bias}:
\begin{itemize}
    \item \textit{Calcium, phosphate, bicarbonate, serum albumin} (8-variable scenario only, routinely drawn on the same chemistry panel as creatinine): nearest value within $\pm$24 hours of the anchor. This window follows APACHE II's precedent of treating all labs drawn within a patient's first 24 hours — including serum creatinine — as one physiologic snapshot \citep{knaus1985apacheii}; no directly analogous 24h-snapshot convention was found for other candidate windows (e.g., KDIGO's 48h AKI change-detection window measures a \emph{trend}, not concurrency, and is not equivalent).
    \item \textit{uACR} (a separately-ordered urine test): nearest value across the patient's entire history. We initially bounded this match to the same hospital admission, following the resolution Tangri et al.\ document for one of their own real-world validation cohorts (\textit{"Geisinger: covariates obtained most closely to index date within a past year were included in models"} \citep{tangri2016multinational}), but this bound caused 96.5\% patient attrition (34{,}332 $\to$ 1{,}220 for the 4-feature scenario) and skewed the outcome rate from 91.9\% to 98.2\% ESRD-positive by disproportionately dropping ESRD-negative patients. We loosened the bound to the patient's whole history, increasing cohort size at the cost of temporal alignment. Because matching is not restricted to measurements before the anchor, these retrospective inputs can include future information relative to that anchor.
    \item Any row missing a required lab within its window is \emph{dropped}, not imputed or flagged — four\_features/eight\_features are complete-case by construction, matching KFRE's own derivation cohort's eligibility criteria \citep{tangri2011predicting}.
\end{itemize}
One row is emitted per qualifying creatinine draw, so a patient monitored multiple times within one admission correctly yields multiple rows, each independently paired with its own nearest in-window value of every other required lab (a sparser lab such as uACR is therefore legitimately reused, at its nearest value, across several rows for the same patient — expected under the landmarking framework above, not a bug).

Table~\ref{tab:cohort_flow} follows the two-column cohort-flow reporting convention of Major et al.\ \citep{major2019kfre_external_validation} and STROBE item 13(a) \citep{vonelm2007strobe}. All three scenarios are event-rich by construction (83--92\% patient-level ESRD-positive), because the source population is itself drawn from a CKD/ESRD diagnosis-code criterion.

\textbf{Cohort size.} The eight-feature scenario is the smallest, with 1{,}213 patients and 1{,}032 observed events. These counts describe the full extracted cohort, before its training/test split, and do not establish the adequacy of the training sample for every architecture. Complete-case selection also changes cohort composition, so cross-scenario differences cannot be attributed to feature count alone.

\section{KFRE Coefficients: Full 8-Variable Form} \label{app:kfre}
\begin{align}
L_8 &= -0.1992\left(\tfrac{\text{age}}{10}-7.036\right) + 0.1602(\text{male}-0.5642) \nonumber \\
&\quad -0.4919\left(\tfrac{\text{eGFR}}{5}-7.222\right) + 0.3364\left(\ln(\text{uACR})-5.137\right) \nonumber \\
&\quad -0.3441(\text{albumin}-3.997) + 0.2604(\text{phosphate}-3.916) \nonumber \\
&\quad -0.07354(\text{bicarb.}-25.57) -0.2228(\text{calcium}-9.355)
\end{align}
Coefficients from \citet{tangri2011predicting}; original-equation baseline survival constants from \citet{tangri2016multinational}: $S_0$(2yr)$=0.9780$ and $S_0$(5yr)$=0.9301$ for eight variables, compared with $0.9750$ and $0.9240$ for four variables. KFRE has no published equation for the twenty\_features scenario and is therefore absent from it.

\section{Benchmarked Models: Full Architectural Detail} \label{app:models}
The neural models, GBSA, and SRF are tuned per scenario (below), so we describe each model's structural form and loss precisely, verified directly against the model implementations, and give the tuned hyperparameter range rather than a single fixed value.

\subsection{Traditional Survival Analysis Models}
Both use \texttt{lifelines} \citep{davidson2019lifelines}; Cox is fitted to counting-process rows and Weibull AFT to each patient's final row.

\textbf{Cox proportional hazards} \citep{cox1972regression} models the hazard of subject $i$ at time $t$ as $\lambda_i(t) = \lambda_0(t)\exp(\boldsymbol{x}_i(t)^\intercal\boldsymbol{\beta})$, where $\boldsymbol{x}_i(t)$ is subject $i$'s (possibly time-varying) covariate row and $\lambda_0(t)$ an unspecified baseline hazard, fit by maximizing the Efron-corrected partial likelihood over all subjects' \texttt{start}/\texttt{stop} intervals directly (\texttt{lifelines}' \texttt{CoxTimeVaryingFitter}, \texttt{id\_col=subject\_id}) rather than one static row per patient — the model consumes every landmark row (Appendix~\ref{app:cohort}), not just each patient's last. We use an L2 penalizer of 1.0 to stabilize the fit against the strongly collinear rows a single patient contributes.

\textbf{Weibull Accelerated Failure Time (AFT)} \citep{anderson1991nonproportional_weibull} instead parameterizes the full survival distribution, $S(t\mid\boldsymbol{x}) = \exp(-(t/\lambda(\boldsymbol{x}))^{\rho})$, $\lambda(\boldsymbol{x}) = \exp(\boldsymbol{x}^\intercal\boldsymbol{\beta})$, fit on the flattened one-row-per-patient (last-observation) frame via \texttt{lifelines}' \texttt{WeibullAFTFitter}. For twenty\_features (whose 40, mostly-zero-filled columns produce a near-singular design matrix), we add an L2 penalizer of 0.1 on the scale parameter's coefficients, the library's own documented remedy; four\_features/eight\_features use the unpenalized default.

\subsection{Deep Survival Analysis Models}
DeepSurv and LogisticHazard use each patient's final row. Dynamic-DeepHit and Hazard Transformer consume full patient histories. RNN-Surv is trained on individual rows represented as length-one sequences and is evaluated on each patient's final row; its recurrent architecture therefore does not model longitudinal histories in this implementation.

\textbf{DeepSurv} \citep{katzman2018deepsurv} is a feedforward network — 1 to 20 \texttt{Linear}$\to$\texttt{ReLU}$\to$\texttt{Dropout} blocks (width 16--256, per-layer dropout 0.1--0.5, all tuned) followed by a single linear output — trained with the Cox partial log-likelihood on its scalar output, so its raw output is already a "higher = riskier" hazard score by construction.

\textbf{Dynamic-DeepHit} \citep{lee2020dynamicdeephit,lee2018deephit} feeds each patient's full covariate row sequence with a padding mask directly into a 2-layer \emph{bidirectional} LSTM (hidden width 16--128, dropout 0.1--0.5, tuned). The concatenated forward/backward outputs at each timestep pass through a \texttt{Linear}$\to$\texttt{Tanh} compression, then a masked additive (Bahdanau-style) attention layer pools the sequence into one context vector per patient, letting the model weight informative visits over uninformative ones. The pooled vector feeds 0--3 further cause-specific \texttt{Linear}$\to$\texttt{ReLU}$\to$\texttt{Dropout} layers (width 16--128, tuned) and a final per-risk linear head producing a probability mass function (PMF) over 5{,}475 daily bins (15 years). Following DeepHit's original formulation \citep{lee2018deephit}, training minimizes a composite loss $L=\alpha L_1+(1-\alpha)L_2$: $L_1$ is the negative log-likelihood of the discretized, right-censored outcome, and $L_2$ is a pairwise ranking loss rewarding the model for assigning an event subject higher cumulative incidence, at that subject's own event time, than any subject not yet failed by then — with $\alpha\in[0.1,1.0]$ tuned per scenario.

\textbf{RNN-Surv} \citep{giunchiglia2018rnn_surv} embeds each timestep through 1--3 linear embedding layers with intervening \texttt{ReLU} activations (width 32--128, tuned) into a 1--3-layer \emph{LSTM} (hidden width 64--256, tuned). A final linear layer maps the last hidden state to logits over $K\in[10,50]$ (tuned) discrete intervals spanning the maximum training follow-up, softmax-normalized into a PMF. As with Dynamic-DeepHit, the loss combines a discrete-time likelihood (observed-interval event probability for events; survival-beyond-interval probability for censored subjects) with a pairwise ranking term over the cumulative incidence function, weighted by a tunable coefficient in $[0.1,0.9]$.

\textbf{LogisticHazard} discretizes training follow-up into 50 bins (\texttt{pycox}'s \texttt{LabTransDiscreteTime}) and learns the log-odds of failure in each bin, via \texttt{pycox}'s off-the-shelf \texttt{MLPVanilla} network (1--3 hidden layers, width $\in\{16,32,64,128\}$, dropout 0--0.5, all tuned, batch size $\in\{32,64,128,256\}$) wrapping its \texttt{LogisticHazard} discrete-time hazard head and cross-entropy-style loss.

\textbf{Hazard Transformer} \citep{hazard_transformer,vaswani2017attention} departs from a sequence-attention reading of "Transformer": each patient's rows are linearly embedded and \emph{mean-pooled} (masked over observed rows) into a single patient representation \emph{before} any attention is applied. This pooled vector is broadcast across 100 fixed query time points spanning $[0,730]$ days, each combined with a learned time-encoding and a sinusoidal positional encoding; 2--6 (tuned) Transformer encoder layers (1--8 heads, tuned) then self-attend \emph{across these 100 time queries}, not across the patient's visits. A per-risk decoder produces a softmax-normalized PMF over the 100 bins (an earlier per-bin independent-sigmoid parametrization collapsed to a constant prediction and was replaced by this bounded-PMF form, the same fix Dynamic-DeepHit's risk heads use). Training minimizes the discrete-time event/censoring negative log-likelihood without an additional ranking term.

\subsection{Ensemble-Learning-Based Models}
All three are \texttt{scikit-survival} \citep{polsterl2020sksurv} estimators fit on the flattened last-observation frame.

\textbf{Gradient Boosting Survival Analysis (GBSA)} \citep{chen2013gradient_GBSA} fits an additive ensemble of regression trees against the Cox partial-likelihood gradient; grid-searched over \texttt{n\_estimators}$\in\{50,100,200,300\}$, \texttt{max\_depth}$\in\{3,5,10,15\}$, \texttt{learning\_rate}$\in\{0.01,0.1,0.2\}$ by cross-validated C-index.

\textbf{Random Survival Forest (SRF)} \citep{pickett2021random_survival_forests,ishwaran2008random} grows an ensemble of survival trees splitting on the log-rank statistic; grid-searched over \texttt{n\_estimators}$\in\{50,100,150\}$, \texttt{max\_depth}$\in\{\text{None},5,10\}$. Since \texttt{scikit-survival} does not expose impurity-based feature importance for this estimator (known to be biased for survival forests), we compute permutation importance for it directly.

\textbf{FastSurvivalSVM} learns a linear ranking function over the covariate space, optimizing a convex upper bound on the concordance index subject to margin and $\ell_2$ constraints. Fit with \texttt{scikit-survival}'s defaults ($\alpha=0.01$), as with the fixed-parameter Cox and Weibull AFT baselines.

\subsection{Model Training and Optimization}
The neural models use 10-trial Optuna searches per scenario, maximizing the implemented C-index objective. DeepSurv, Dynamic-DeepHit, RNN-Surv, and Hazard Transformer select trials on training-set C-index and stop when training loss fails to improve (patience five epochs, maximum 50 epochs). LogisticHazard is fitted for 100 epochs and uses the designated test partition both as \texttt{val\_data} and for trial/checkpoint selection. Its reported test performance therefore has model-selection leakage and is descriptive rather than an unbiased held-out estimate. GBSA and SRF use five-fold grid search within the training partition, selected by C-index. FastSurvivalSVM, Cox, and Weibull AFT use fixed hyperparameters. All neural models use Adam.

\section{Evaluation Definitions and Scope} \label{app:evaluation}
\textbf{Repetitions and aggregation.} We report all three scenarios from the completed clinical-validity analyses for repetitions 1--5. Each scenario uses an 80/20 split by patient ID with \texttt{random\_state=42}; the repetitions reuse the same partition. Results are the arithmetic mean and sample standard deviation ($n=5$, denominator $n-1$) of the five report values. They describe run-to-run variation on a fixed partition, not uncertainty across independent cohorts or externally validated confidence intervals.

\textbf{Prediction unit.} Every model contributes one test prediction per patient. Snapshot models use the final available row; Dynamic-DeepHit and Hazard Transformer use the available row history and the outcome recorded on its final row. RNN-Surv uses the final row as a length-one sequence. Thus the scoring unit is shared, while input histories differ. Durations remain measured from the extraction's first-lab anchor rather than being reset at a prospective prediction landmark. This retrospective evaluation does not establish prospective two-year prediction from a common baseline.

\textbf{Metrics.} Harrell's C-index is calculated over each model's available patient follow-up, without truncation at 730 days, after orienting scalar scores so higher values indicate greater risk. Mean cumulative/dynamic AUC uses that scalar score on a daily grid from day 1 through day 728 (or the shorter observed support), with the survival-weighted mean returned by \texttt{scikit-survival}. Score definitions remain model-specific: for example, Dynamic-DeepHit uses day-730 cumulative incidence, Hazard Transformer uses day-365 cumulative incidence, RNN-Surv uses expected event earliness, and LogisticHazard evaluates survival near the scored partition's median duration.

The integrated Brier score (IBS) uses 50 equally spaced times from day 1 to day 730. For every model, including KFRE, the reported IBS uses a score-derived survival curve $\hat S(t\mid r)=\exp[-r^*\hat H_0(t)]$, where $r^*=\max(r-\min(r_{\mathrm{train}}),10^{-6})$ and $\hat H_0$ is a Breslow-style cumulative baseline hazard fitted to that model's training-patient scores and outcomes. These IBS values evaluate this additional conversion rather than each model's native survival probabilities. For IBS only, test follow-up beyond the maximum training duration is administratively censored there. The inverse-probability-of-censoring reference for IBS and AUC is constructed from raw training rows, while test predictions are per patient; this mismatch remains a limitation of the completed analysis.

\textbf{Calibration and exploratory decision curves.} Calibration compares mean predicted risk with Kaplan--Meier-observed risk within risk deciles at 730 and 1{,}825 days. It uses native horizon probabilities where supported, with the score-derived survival conversion above as fallback. Unlike the IBS comparison, native KFRE calibration therefore uses the published equation directly. The generated decision-curve analyses \citep{vickers2006decisioncurve} include treat-all, treat-none, and eGFR referral rules ($<$30 or $<$45), but model curves use patient-level predictions whereas treat-all and eGFR comparators use raw test rows. The eGFR calculation also uses its referral fraction as an implied risk threshold. These unmatched units and thresholds preclude clinical-utility comparisons; the generated curves are exploratory and are not used to infer superiority over referral baselines. A competing-risk analysis for death before ESRD was not performed because exported durations omit the absolute calendar anchor needed to align death dates.
